\section*{結言}
今回の10件では、モデル自体を操作する研究と、エージェントが長期的に働くための開発基盤という二つの流れが交差した。Weight-Space Programmingの報告はモデル編集という新しい方向性を示し、ML4の技術レポートはオープンウェイトモデルの評価材料を増やした。ただし、注目を集めた実験値や次世代GPTの公開時期には、独立再現や一次確認が残っている。

Beacon、Spec Kit、Creator Mode、OpenCode、TermiX、WonderDriveに共通するのは、単発の生成結果よりも作業の継続性、制御、拡張、共有を重視する姿勢である。一方、AI学習オプトアウト信号の議論は、技術的な合意だけでは利用者の選択が直ちに実現するとは限らないことを示した。

Grokは今回の収集範囲を\textit{Good}と評価したが、X検索の完全性には限界がある。未確認の製品仕様、性能値、企業発言は確定情報と混同せず、今後の公式資料や実装で確かめる必要がある。
